\section{The finite-torus and whole-plane bounds}\label{sec:conclusion}

We now apply the geometric comparison to the certified scalar bounds,
then pass from square tori to the plane. The square-torus conclusion of
Proposition~\ref{dual:assembly} already incorporates the separation of
minimizers and the passage to arbitrary configurations.

\begin{proof}[Proof of Theorem~\ref{thm:torus}]
The radial data of Section~\ref{sec:calibration} satisfy
\eqref{dual:radial-origin}--\eqref{dual:radial-sum} with lower distance
$d_0=b=28209/100000$. Theorem~\ref{comp:scalar} gives $K<0$ and both
bounds in \eqref{cal:scalar-bounds}. By
Lemma~\ref{dual:signed-endpoints} at $d_0=b$ and the primitive identity
\eqref{cal:sector-primitive}, these bounds give
\eqref{dual:sector-input} for every sector with $d\geq b$, including
the signed endpoint cases.

Since $b<1/(2\sqrt\pi)=d_*$, as follows from $\pi<3.1416$, the
square-torus conclusion of Proposition~\ref{dual:assembly} applies.
For every $n\geq2$ and every configuration of $n$ distinct points it
gives, using the exact triangular calibration
\eqref{cal:triangular-energy},
\[
\mathcal W_n(h)\geq n(6K+\lambda+2\pi\mu)=nW(E_\triangle).
\]
No separation hypothesis is imposed on this configuration. The same
proposition gives the energy penalty $n|c|^2/2$ for a constant harmonic
component and preserves the bound under a fixed orthogonal rotation.
\end{proof}

\begin{proof}[Proof of Theorem~\ref{thm:main}]
The periodic triangular field belongs to the admissible class and has
finite energy by Lemmas~\ref{red:local} and \ref{red:periodic-average}.
Thus $C=W(E_\triangle)$ is a finite real number.  Apply
Proposition~\ref{red:reduction} to Theorem~\ref{thm:torus} with this
$C$. It gives the lower bound for every field
in the stated admissible class and for each fixed cutoff family
\eqref{eq:cutoffs}. The triangular field attains the lower bound.
\end{proof}

The comparison establishes the minimum of this logarithmic
renormalized energy. It does not classify all configurations
having that energy per unit area. In particular it supplies no
uniqueness or defect-exclusion assertion.

\subsection{The spherical logarithmic-energy constant}

\begin{proof}[Proof of Corollary~\ref{cor:bhs}]
Write $\mathcal A_m^{\mathrm{BS}}$ and $W_{\mathrm{ball}}$ for the
admissible class and ball-cutoff energy of B\'etermin and
Sandier~\cite[Definitions~2.1--2.4]{BeterminSandier2018}. Their fields
satisfy $\operatorname{div}E=2\pi(\nu_\Lambda-m)$ and $\curl E=0$,
with the same charge-growth condition and local puncture functional as
here. At $m=1$ the admissible class is our $\mathcal A_1$; the local
integrability conventions agree by Lemma~\ref{red:local}. The outer
average, however, is over centered balls rather than squares.

The component rotation $j\mapsto F=(j_2,-j_1)$ is a bijection from
$\mathcal C_1$ to $\mathcal A_{1/(2\pi)}^{\mathrm{BS}}$ and preserves
the local energy.
Thus \eqref{red:M} identifies $M$ with
$\min_{\mathcal A_{1/(2\pi)}^{\mathrm{BS}}}W_{\mathrm{ball}}$.
The scaling formula \cite[equation~(2.4)]{BeterminSandier2018} gives
\begin{equation}\label{eq:bhs-minimum-normalization}
\min_{\mathcal A_1^{\mathrm{BS}}}W_{\mathrm{ball}}
=2\pi\left(M-\frac14\log(2\pi)\right).
\end{equation}
Lemma~\ref{red:plane-bound} bounds every square-cutoff energy $W(E)$
below by the right-hand side. Conversely, periodic averaging applied
to the scaled square-periodic sequence in the proof of
Proposition~\ref{red:reduction} shows that its square-cutoff energies
converge to that value. Theorem~\ref{thm:main} therefore identifies it
with $W(E_\triangle)$. This uses the common ball and square minimum in
\eqref{red:M}, not equality of those two energies for an arbitrary
field. Lemma~\ref{red:periodic-average} does give equality for the
periodic field $E_\triangle$, so this covolume-one triangular field
attains the minimum in \eqref{eq:bhs-minimum-normalization}. Its value
also agrees with the periodic lattice notation in B\'etermin and
Sandier: replacing their gradient $\nabla H_\Lambda$ by the admissible
field $-\nabla H_\Lambda$ leaves the local squared-norm functional
unchanged.

Theorem~1.5 of \cite{BeterminSandier2018} gives the form of
\eqref{eq:bhs-asymptotic} with $C$ in place of $C_{\mathrm{BHS}}$, where
\[
C=\frac1\pi\min_{\mathcal A_1^{\mathrm{BS}}}W_{\mathrm{ball}}
       +\frac12\log\pi+\log2,
\]
and states that $C=C_{\mathrm{BHS}}$ exactly when the density-one
triangular lattice attains this minimum. Its energy is precisely the
ordered-pair sum on the unit sphere in $\R^3$ used here. In particular,
the $-n(n-1)\log2$ term in its passage from the radius-$1/2$ sphere to
the unit sphere contributes both $-n^2\log2$ and $+n\log2$; no further
pair-counting or radius factor is needed. The equality clause now gives
\eqref{eq:bhs-constant} and proves the corollary.
\end{proof}
